Let $x\in G(\A)^1$, let $c\ge0$, and let $\delta_0\in G(\Q)$ satisfy $\alpha(\HB(\delta_0x))=\operatorname{ht}(x)$ (such $\delta_0$ exist by
\texttt{I.red.thm:height-attained}). Then:
\begin{enumerate}
  \item[(a)] if $\operatorname{ht}(x)>c$, the cosets $B(\Q)\delta\in B(\Q)\backslash G(\Q)$ with $\alpha(\HB(\delta x))>c$ are exactly the one coset
  $B(\Q)\delta_0$; that is, for $\delta\in G(\Q)$, $\alpha(\HB(\delta x))>c$ if and only if $\delta\delta_0^{-1}\in B(\Q)$;
  \item[(b)] if $\operatorname{ht}(x)\le c$, then $\alpha(\HB(\delta x))\le c$ for every $\delta\in G(\Q)$.
\end{enumerate}
In particular, at most one coset $B(\Q)\delta$ satisfies $\alpha(\HB(\delta x))>c$, and one does if and only if $\operatorname{ht}(x)>c$.
